\subsection{Persisting Relations} \label{s:pr}
Once we've done the hard work of extracting key \md{} and presenting mechanisms to the user (both human and application) to define relations, we need to select a means of persisting these relations in a manner that makes sense for the data type. In our model we have chosen a \gdb{} for this function as the graph model provides an excellent means of describing relation. Graph database models can be defined as those in which data structures for the schema and instances are modeled as graphs or generalizations of them, and data manipulation is expressed by graph-oriented operations and type constructors \cite{Angles:2008:SGD:1322432.1322433}.

Graph databases are described as ``whiteboard friendly.'' Drawing circles and connecting them with lines on a whiteboard is how we visualize graphs. We denote data points as \emph{nodes} and connect (relate) nodes with \emph{links}\footnote{Graph theory is a branch of discrete mathematics. Nodes are referred to as \emph{vertices} and what we call links are referred to as \emph{edges}, where edges connect pairs of distinct vertices. A simple graph (such as we'll be using in our descriptions) with $V$ vertices has at most $V(V-1)/2$ edges. We won't be dealing with graphs mathematically other than to note that a graph is defined by its vertices and its edges, not by the way we choose to draw it \cite{DBLP:books/daglib/0001826}.}, as illustrated in Figure \ref{fig:epi} (p. \pageref{fig:epi}). Graph databases, such as NEO4j \cite{Webber:2012:PIN:2384716.2384777}, are well suited to \ud{} as they are typeless and have no set schema, while providing several advantages for our application.

First, the graph data model is useful when the interconnectivity of data and the ability to discover the \emph{relationships between values} is important, rather than simply commonality \emph{among value sets} typical to relational models. Discovering relations is fundamental to our design, so its important we select a data structure that can readily describe relations with good performance. Unlike join operations in relational databases or map-reduce operations in other databases, graph traversals are in constant time \cite{redmond2012seven}. Graph databases make it easy to discover centrality, where we measure individual nodes against a full graph. (The most famous centrality algorithm is Google PageRank \cite{ilprints422}. In \S\ref{s:models} we use centrality in an epidemiology example.) Finally, graph databases have been shown to perform full-text character searches significantly faster than relational databases \cite{Vicknair:2010:CGD:1900008.1900067}.

Popular examples of publicly available \apps{} that make use of graph databases include Freebase\footnote{\href{http://www.freebase.com}{http://www.freebase.com}} and the Disease Ontology database\footnote{\href{http://disease-ontology.org}{http://disease-ontology.org}}. Freebase, created by Metaweb and acquired by Google in 2010, is a large community \md{} \db{} described as ``an open shared database of the world's knowledge.'' Disease Ontology, created by the University of Maryland School of Medicine, represents a comprehensive knowledge base of 8043 inherited, developmental and acquired human diseases \cite{schriml2012disease}.
